% !TEX root = ../Hamiltonian_of_mean_force.tex
\section{The Hamiltonian of Mean Force\label{sec:qubit_closure}}
\subsection{Closure as an operator-algebra statement\label{subsec:closure_algebra}}

With this, we may now give the general expression for the Hamiltonian of Mean-Force. 
\begin{equation}
    -\beta H_{\mathrm{MF}}(\beta)
    = \log\!\big(e^{-\beta H_Q}\,e^{\Delta(\beta)}\big).
    \label{eq:HMF_log_def}
\end{equation}



If $H_Q$ and $\Delta$ commuted, the logarithm would trivially give $-\beta H_Q+\Delta$. In general
they do not, and the Baker--Campbell--Hausdorff (BCH) expansion
\begin{equation}
    \log(e^A e^B) = A + \Phi(\mathrm{ad}_A)\,B + O(B^2),
    \qquad
    \Phi(x)\equiv\frac{x}{1-e^{-x}},
    \label{eq:BCH_linear}
\end{equation}
generates a tower of commutators of $A=-\beta H_Q$ acting on $B=\Delta$. To understand the structure of these corrections, we invoke the partition of the influence exponent into its secular and mixing parts, $\Delta = \Delta_{\parallel} + \Delta_{\perp}$. The structural partitioning of these mean-force corrections is illustrated in Fig.~\ref{fig:hmf_flowchart}. Because the secular piece $\Delta_{\parallel}$ strictly commutes with $A$, its adjoint action is annihilated ($\mathrm{ad}_A(\Delta_{\parallel}) = 0$).

Consequently, the generating function evaluated at $x=0$ yields $\Phi(0)=1$, allowing the parallel component to push straight through the linear BCH expansion:
\begin{equation}
    -\beta H_{\mathrm{MF}} \approx \big(-\beta H_Q + \Delta_{\parallel}\big) + \Phi(\mathrm{ad}_{-\beta H_Q})\,\Delta_{\perp} + O(\Delta^2).
    \label{eq:BCH_split_linear}
\end{equation}
The secular zeroth-order shift $\Delta_{\parallel}$ thus groups directly with the bare system (reminiscent of the weak-coupling Lamb shift). Crucially however, while it acts geometrically like a classical energy shift, this term is still fundamentally \textit{quantum} in origin because it is weighted by the bath Green functions, dynamically entangling the bare spectrum with the full spectral density of the environment. We can understand this directly by reference to Eq.~\eqref{eq:Delta_selfenergy}. By defining a bath-dressed coupling operator $\tilde f$ with matrix elements $\tilde f_{ij} = f_{ij}\sqrt{R(\omega_{ij})}$, the secular exponent takes the exact quadratic form $\Delta_{\parallel} = (\tilde f \tilde f^\dagger)_{\text{diag}}$. In the macroscopic limit where the system commutes with its coupling ($[H_Q, f]=0$), the matrix $f$ is entirely diagonal in the energy basis ($\omega_{ij}=0$), and this dressed square collapses precisely to $\Delta_{\parallel} = \beta\lambda f^2$, where $R(0) = \beta\lambda$ is the explicit reorganisation energy defined in Eq.~\eqref{eq:int_K_2lambda}. The zeroth order BCH term therefore exactly recovers the classical commuting shift $H_{\mathrm{MF}} = H_Q - \lambda f^2$ of Eq.~\eqref{eq:Heff_commuting_limit}, and at zeroth order the non-commuting correction can be obtained by replacing $f \to \tilde f$. The remaining coherence-generating effects are then driven exclusively by the action of $\Phi$ on the perpendicular, off-diagonal open-chain fragments $\Delta_{\perp}$.



There is a saitsfyingly clean expression for these non-commuting corrections in the system energy eigenbasis. The off-diagonal basis operators $\ket{i}\!\bra{j}$ are exact eigenoperators of the adjoint map: $\mathrm{ad}_{-\beta H_Q}\ket{i}\!\bra{j} = -\beta [H_Q, \ket{i}\!\bra{j}] = -\beta\omega_{ij}\ket{i}\!\bra{j}$. Consequently, the entire infinite tower of nested commutators generated by $\Phi(\mathrm{ad}_{-\beta H_Q})$ trivially evaluates to a simple scalar multiplication for each matrix element of the perpendicular fragment:
\begin{equation}
    \Phi(\mathrm{ad}_{-\beta H_Q})\Delta_{\perp} = \sum_{i\neq j} \Phi(-\beta\omega_{ij})\,\Delta_{ij}\ket{i}\!\bra{j},
\end{equation}
where $\Phi(x) = x/(1-e^{-x})$. The off-diagonal elements of the linear-in-$\Delta$ effective Hamiltonian are therefore simply the bare bath-induced transitions $\Delta_{ij}$, parametrically re-weighted by the thermal asymmetry factor $\Phi(-\beta\omega_{ij}) = -\beta\omega_{ij}/(1-e^{\beta\omega_{ij}})$.

This scalar evaluation exhibits a manifestly physical symmetry. If we reconstruct the effective matrix element using the symmetrised influence operator $S_{ij} = e^{-\beta\omega_{ij}/2}\Delta_{ij}$, the thermal asymmetry factor exactly absorbs the detailed-balance twist. The (off-diagonal) components on the effective Hamiltonian are therefore simply the bare bath-induced transitions $\Delta_{ij}$, parametrically re-weighted by the thermal asymmetry factor $\Phi(-\beta\omega_{ij}) = -\beta\omega_{ij}/(1-e^{\beta\omega_{ij}})$:
\begin{equation}
\begin{split}
    (H_{\mathrm{MF}})_{ij} &= -\frac{1}{\beta}\Phi(-\beta\omega_{ij}) e^{\beta\omega_{ij}/2} S_{ij} \\
    &= -\frac{\omega_{ij}/2}{\sinh(\beta\omega_{ij}/2)}\, e^{-\beta\omega_{ij}/2} \Delta_{ij} .
\end{split}
    \label{eq:HMF_linear_off_diag}
\end{equation}
Because $x/\sinh(x)$ is an explicitly even geometric function of the energy gap $\omega_{ij}$, Eq.~\eqref{eq:HMF_linear_off_diag} guarantees that the linearly truncated effective Hamiltonian is perfectly Hermitian. Moreover, this exact thermal weighting reveals that bath-induced coherent mixing between levels is heavily exponentially suppressed for high-frequency transitions ($|\beta\omega_{ij}| \gg 1$), dynamically localising the leading mean force corrections to the thermodynamic energy shell.

This fortuitous cancellation reveals a deep structural connection between quantum thermodynamics and Lie algebra. The fundamental detailed balance condition of the bath imposes a thermal asymmetry $e^{-\beta\omega}$ on the bare transition amplitudes. Conversely, the BCH expansion---a purely abstract property of non-commuting logarithms---generates the parametric weighting $\Phi(-\beta\omega) \propto (1-e^{\beta\omega})^{-1}$. The fact that the Lie-algebraic geometry of the thermal exponential perfectly mirrors - and thereby absorbs - the thermodynamic detailed-balance twist of the environment is striking, and illustrative of the high degree of mathematical harmonisation underlying the mean-force formulation.

\begin{figure}[t]
\centering
\begin{tikzpicture}[
    node distance=1.0cm and 0.4cm,
    font=\footnotesize,
    box/.style={rectangle, draw, rounded corners, minimum width=3.8cm, minimum height=1.3cm, text centered, text width=3.4cm, fill=white, thick},
    titlebox/.style={rectangle, draw=blue!60, fill=blue!5, rounded corners, minimum width=8.0cm, minimum height=1.1cm, text centered, text width=7.4cm, thick},
    arrow/.style={-Latex, thick, shorten >=2pt, shorten <=2pt},
    groupbox/.style={rectangle, draw, rounded corners, dashed, inner sep=15pt, fill=gray!5, thick},
    highlight/.style={fill=red!10, draw=red!60},
    glow/.style={fill=yellow!10, draw=orange!60},
    scale=0.85, transform shape
]

% Top level input
\node (delta) [titlebox, minimum height=0.8cm] {\textbf{Bare Transition Amplitude Matrix} $\Delta$};

% Splitting
\node (parallel) [box, below=3.2cm of delta, xshift=-2.3cm, glow] {\textbf{Secular Sector} ($\Delta_{\parallel}$) \\ \vspace{2pt} \textit{Diagonal resonant} ($i = j$) \\ $\Delta_{ii} = \beta\lambda f_{ii}^2$};
\node (perp) [box, below=3.2cm of delta, xshift=2.3cm, highlight] {\textbf{Mixing Sector} ($\Delta_{\perp}$) \\ \vspace{2pt} \textit{Off-diagonal non-resonant} ($i \neq j$) \\ $\Delta_{ij} \propto \lvert f_{ij}\rvert^2 R(\omega_{ij})$};

% Processes
\node (bch1) [box, below=3.0cm of parallel, fill=yellow!5, draw=orange!60] {\textbf{Macroscopic Ext.} \\ $\Phi(\mathrm{ad}_{H_Q}) = 1$ \vspace{2pt} \\ Classical shift };
\node (bch2) [box, below=3.0cm of perp, fill=red!5, draw=red!60] {\textbf{Lie Generator} \\ $\Phi(\mathrm{ad}_{-\beta H_Q}) \to \frac{\omega_{ij}/2}{\sinh(\beta\omega_{ij}/2)}$ \vspace{2pt} \\ Absorbs detailed balance};

% Output
\node (hmf) [titlebox, bottom color=blue!10, top color=white, below=3.2cm of $(bch1.south)!0.5!(bch2.south)$] {\textbf{Effective Mean-Force Hamiltonian} $H_{\mathrm{MF}}$ \\ \vspace{4pt} 
Diag: $E_i - \lambda f_{ii}^2$ \qquad Mix: $-\frac{\omega_{ij}/2}{\sinh(\beta\omega_{ij}/2)} \, e^{-\beta\omega_{ij}/2}\Delta_{ij}$ };

% Arrows
\draw [arrow] (delta.south) -- ++(0,-1.5) coordinate (branch);
\draw [arrow] (branch) -| (parallel.north);
\draw [arrow] (branch) -| (perp.north);

\path (branch) -- (branch -| parallel.north) node[midway, above, font=\scriptsize\itshape] {Diagonal proj.};
\path (branch) -- (branch -| perp.north) node[midway, above, font=\scriptsize\itshape] {Off-diagonal proj.};

\draw [arrow] (parallel) -- (bch1);
\draw [arrow] (perp) -- (bch2);

\draw [arrow] (bch1.south) -- ++(0,-1.5) -| (hmf.north);
\draw [arrow] (bch2.south) -- ++(0,-1.5) -| (hmf.north);

% Background labels
\begin{scope}[on background layer]
    \node (bg1) [groupbox, fill=gray!10, draw=gray!50, fit=(parallel) (perp)] {};
    \node [above, font=\bfseries\color{gray!70}, inner sep=3pt] at (bg1.north) {Influence Partitioning};
    
    \node (bg2) [groupbox, fill=blue!5, draw=blue!40, fit=(bch1) (bch2)] {};
    \node [above, font=\bfseries\color{blue!50}, inner sep=3pt] at (bg2.north) {BCH Expansion Action};
\end{scope}

\end{tikzpicture}
\caption{Structural decomposition of the linear-order Mean-Force Hamiltonian. The parallel (secular) sector generates classical Lamb-like commuting phase shifts, entirely isolated from detailed balance effects. Concurrently, the perpendicular (mixing) sector passes through the Baker--Campbell--Hausdorff non-commutative generator, producing a geometric Lie scaling that perfectly absorbs the detailed balance thermal twist.}
\label{fig:hmf_flowchart}
\end{figure}

The foregoing parametric logic naturally generalises directly to all higher orders of the BCH expansion. Because the basis operators $\ket{i}\!\bra{j}$ are strict eigenoperators of the adjoint map, the evaluation of any highly nested Lie commutator string spanning $O(\Delta^n)$ simply produces ordinary matrix products of the abstract fragments $\Delta_{ik_1} \Delta_{k_1 k_2} \dots \Delta_{k_{n-1}j}$, completely bypassing the need for functional derivatives. Over the energy eigenbasis, the entire non-linear algebraic sum reduces to evaluating $n$-step discrete pathways through the bare spectrum, where each composite transition amplitude is identically parametrically weighted by a universal, symmetric multivariate generalisation of $\Phi$ determined solely by Lie combinatorics.  

While this geometric pathing approach is exceptionally suggestive---and likely a numerically robust framework for resolving perturbative components---the abstract operator structure is more directly revealing of the core algebraic symmetries. Because the adjoint action obeys the rule $\mathrm{ad}_{H_Q}(f_n f_m) = f_{n+1}f_m + f_n f_{m+1}$ (where $f_n \equiv \mathrm{ad}_{H_Q}^n f$), every nested BCH commutator stays within the algebra generated by the adjoint chain of the coupling. This ensures that even high-order nonlinear corrections introduce \emph{no new bath objects}: the environment continues to enter exclusively through the fundamental spectral functions $\mathcal{K}(\omega)$ and $R(\omega)$, with coefficients determined solely by universal Lie combinatorics.

The closure content of $H_{\mathrm{MF}}(\beta)$ is therefore the following: the entire operator
lives in the finite-dimensional algebra generated by $\{H_Q,\Delta\}$, with $\Delta$ quadratic
in $f$ and the bath entering exclusively through the two spectral objects $\mathcal{K}(\omega)$
and $R(\omega)$. Whether this algebra is itself finite-dimensional - and hence whether $H_{\mathrm{MF}}$
admits a closed form in a restricted operator family - depends on the structure of $f$ and $H_Q$ This statement leads to an important corollary. For continuous variable systems $H_Q = p^2/2m + V(x)$, if the potential $V(x)$ is a polynomial of degree greater than 2, the adjoint chain $\{f_n\}$ generally does not close in any finite-dimensional Lie algebra; instead, it generates the full universal enveloping algebra. Consequently, no finite representation of the mean force exists within the original operator class. This obstruction is analogous (in spirit) to the Groenewold--Van Hove no-go theorem on consistent quantisation of polynomial observables: the failure to close the Lie algebra under adjoint action corresponds to the impossibility of mapping the quantum evolution of non-quadratic potentials to a finite phase space flow. For such potentials, there is no \emph{local operator} that exactly describes the reduced equilibrium state.

From a practical standpoint, the closure of the adjoint-chain is not necessary for an arbitrarily accurate approximation of $H_{\mathrm{MF}}(\beta)$, as the BCH logarithm $\log(e^{-\beta H_Q}e^\Delta)$ still yields a controlled expansion. From this algebraic perspective, the closure problem offers a unifying view of strong-coupling approximations. Numerical schemes such as the polaron transformation~\cite{weissQuantumDissipativeSystems2012,nazirReactionCoordinateMapping2018}, hierarchical equations of motion~\cite{tanimuraReducedHierarchicalEquations2014}, and chain-mapping/DMRG techniques~\cite{XuVadimovStockburgerAnkerhold2026} can be understood as distinct choices of a truncated operator family $\mathcal{A}$. The polaron frame corresponds to dressing $\mathcal{A}$ with coherent displacements, HEOM truncates the memory depth of the influence functional (dual to the operator degree in the adjoint chain), and chain mappings truncate the spatial extent of the harmonic bath. The `closure' problem is thus equivalent to identifying a low-dimensional subalgebra that approximately contains the logarithm of the quenched propagator.


